% Generated draft: requires mathematical and source review.
\documentclass[11pt]{article}
\usepackage[margin=25mm]{geometry}
\usepackage{fontspec}
\setmainfont{Latin Modern Roman}
\usepackage{amsmath,amssymb,mathtools,unicode-math}
\setmathfont{Latin Modern Math}
\usepackage{enumitem,needspace,xurl,hyperref}
\setlength{\parindent}{0pt}
\setlength{\parskip}{5pt}
\setlength{\emergencystretch}{3em}
\newcommand{\N}{\mathbb{N}}
\newcommand{\Z}{\mathbb{Z}}
\newcommand{\Q}{\mathbb{Q}}
\newcommand{\R}{\mathbb{R}}
\newcommand{\C}{\mathbb{C}}
\newcommand{\F}{\mathbb{F}}
\newcommand{\A}{\mathbb{A}}
\newcommand{\PP}{\mathbb{P}}
\newcommand{\E}{\mathbb{E}}
\newcommand{\eps}{\varepsilon}
\newcommand{\dd}{\,\mathrm{d}}
\begin{document}
\section*{UnsolvedMath ID: 2000017}
\textbf{Source problem code:} catalogue entry
\subsection{Definitions and mathematical statement}
Let \(\psi(A,\theta)\) denote an elastic crystal's free-energy density, where \(A\) is a deformation gradient and \(\theta\) is temperature. The problem asks which boundary conditions lead to attainment of the associated energy minimum and which lead to nonattainment.

To express the target once the model is fixed, let \(\Omega\subset\R^d\) be a reference body, fix \(\theta\), and let \(\mathcal A_B\) be the admissible deformations for boundary data \(B\). In a purely bulk, fixed-temperature model one would consider
\[
I_\theta(y)=\int_\Omega\psi(\nabla y(x),\theta)\,\dd x,
\qquad m_B=\inf_{y\in\mathcal A_B}I_\theta(y).
\]
For nonempty \(\mathcal A_B\) and finite \(m_B\), attainment means
\[
\exists y_B\in\mathcal A_B:\quad I_\theta(y_B)=m_B,
\]
and nonattainment means no such admissible deformation exists. The research task is to characterize the boundary data in these two classes for the intended crystal free energies. The displayed bulk model is conditional notation, not a claim that the excerpt supplies a complete energy functional.

The record does not specify the crystal symmetry/constitutive class, the body, spatial dimension, temperature regime, admissible function space, determinant or injectivity constraints, external loading, or whether boundary data are Dirichlet, traction, or mixed. These choices affect compactness, lower semicontinuity, and the meaning of an admissible minimizer. The task is not a claim that coercivity alone decides attainment.

The supplied moderate-confidence OPEN-TRIAGE summary reports that standard coercive (quasi)convex existence frameworks and some noncoercive mixing examples are understood, but that a systematic characterization covering crystal free energies and all relevant boundary conditions is not established. No particular theorem or example, with its hypotheses, is supplied.
\subsection{Short English statement}
For elastic-crystal free energies, classify boundary data according to whether the infimum is achieved by an admissible deformation. The model and admissible boundary-condition class must be fixed before this becomes a definite classification problem.
\subsection{Sources}
\begin{itemize}
\item Local UnsolvedMath registry, permanent ID \(2000017\), exact source code \texttt{catalogue entry}, ``Some Open Problems in Elasticity --- Elastic-crystal free energies.''
\item Ball, elasticity problem list (2002), Problem 17; \url{https://people.maths.ox.ac.uk/ball/Articles%20in%20Conference%20Proceedings%20and%20Books/JMB%202002%20re%20Marsden%2060th.pdf}. PDF-text extraction; recorded access 2026-07-29.
\item Supplied moderate-confidence OPEN-TRIAGE summary, without specific authoritative references for the mentioned existence frameworks or noncoercive mixing examples.
\end{itemize}
\end{document}
